% ===== Document Preamble =====
\documentclass[11pt]{article}
\usepackage[margin=0.5in]{geometry}
\usepackage{amsmath,amssymb}
\usepackage{xcolor}
\usepackage{enumitem}
\usepackage{needspace}
\usepackage{comment}

% ===== Colors =====
\definecolor{AnsRed}{RGB}{170,0,0}
\definecolor{WrongBlue}{RGB}{0,0,160}
\definecolor{SkipGray}{RGB}{120,120,120}

% ===== Choice list styling (tight, consistent) =====
\setlist[enumerate]{%
  leftmargin=1.5em,
  itemsep=0.25em,
  topsep=0.25em,
  label=\textbf{\Alph*.}
}

% ===== Style Wrappers =====
\newcommand{\QuestionStyle}[1]{{\large\bfseries\boldmath #1}}
\newcommand{\AnswerStyle}[1]{{\large\bfseries\boldmath #1}}
\newcommand{\ReasonStyle}[1]{{\normalsize #1}}

% ===== Show / Hide Reasoning Toggle =====
\newif\ifshowreason
\showreasontrue      % Instructor version
% \showreasonfalse   % Student version

% ===== Choice Item Macros (Answer + Reasoning) =====
% Usage:
%   \Correct{Answer}{Reasoning}
%   \Wrong{Answer}{Reasoning}

\newcommand{\Correct}[2]{%
  \item \AnswerStyle{\textcolor{AnsRed}{#1}}%
  \ifshowreason\\[-0.2em]
    \ReasonStyle{\textcolor{AnsRed}{#2}}%
  \fi
}

\newcommand{\Wrong}[2]{%
  \item \AnswerStyle{\textcolor{WrongBlue}{#1}}%
  \ifshowreason\\[-0.2em]
    \ReasonStyle{\textcolor{WrongBlue}{#2}}%
  \fi
}

% ===== Multiple-Choice Question Block Macro =====
% Usage: \QBlock{Q1}{stem...}{<items>}{answer key}

\newcommand{\QBlock}[4]{%
  \Needspace{8\baselineskip}%
  \noindent
  \begin{minipage}{\linewidth}
    \QuestionStyle{#1.~#2}

    \vspace{0.25em}
    \begin{enumerate}
      #3
    \end{enumerate}

    \vspace{0.1em}
    \noindent\textcolor{AnsRed}{\textit{\textbf{Answer:}~#4}}

    \vspace{0.25em}
    \hrule height 0.4pt
  \end{minipage}
  \vspace{0.75em}
}

% ===== Open-Ended Question Block Macro =====
% Usage: \OEBlock{Q1}{stem...}{answer text}

\newcommand{\OEBlock}[3]{%
  \Needspace{5\baselineskip}%
  \noindent
  \begin{minipage}{\linewidth}
    \QuestionStyle{#1.~#2}

    \vspace{0.25em}
    \noindent\textcolor{AnsRed}{\textit{\textbf{Answer:}~#3}}

    \vspace{0.25em}
    \hrule height 0.4pt
  \end{minipage}
  \vspace{0.75em}
}

% ===== Multi-Part Question Block Macro =====
% Usage: \MPBlock{Q1}{stem...}{<parts>}

\newcommand{\MPBlock}[3]{%
  \Needspace{8\baselineskip}%
  \noindent
  \begin{minipage}{\linewidth}
    \QuestionStyle{#1.~#2}

    \vspace{0.25em}
    #3

    \vspace{0.25em}
    \hrule height 0.4pt
  \end{minipage}
  \vspace{0.75em}
}

% ===== MC Sub-Part Macro =====
% Usage: \MCPart{a}{stem...}{<items>}{answer key}

\newcommand{\MCPart}[4]{%
  \par\vspace{0.35em}
  \noindent\textbf{(#1)}~#2
  \begin{enumerate}
    #3
  \end{enumerate}
  \vspace{0.1em}
  \noindent\textit{\textbf{Answer:}~#4}
  \par
}

% ===== Open-Ended Sub-Part Macro =====
% Usage: \OEPart{a}{stem...}{answer text}

\newcommand{\OEPart}[3]{%
  \par\vspace{0.35em}
  \noindent\textbf{(#1)}~#2
  \par\vspace{0.15em}
  \noindent\textit{\textbf{Answer:}~#3}
  \par
}

% ===== Skipped Question Note Macro =====
% Usage: \SkipBlock{Q1}{reason}

\newcommand{\SkipBlock}[2]{%
  \noindent\textcolor{SkipGray}{\textit{[#1 omitted: #2]}}%
  \vspace{0.5em}\par
}

% ===== Document Start =====

\begin{document}

\begin{center}
{\huge \textbf{Physics Isomorphic Questions}}\\[0.25em]
{\Large Multiple-Choice Question Set}\\[0.25em]
{\large Questions generated below. Multiple-choice answers in RED; incorrect answers in BLUE.}
\end{center}

\vspace{0.75em}

\QBlock{Q1}{Which of the following statements describes a reasonable physics idealization?}{%
\Wrong{A bicycle wheel can be treated as a point mass when analyzing how different points on the rim move during rotation.}{This ignores the extended shape of the wheel, even though the motion of different parts of the rim is exactly what matters in rotational analysis.}
\Correct{A train traveling $120\,\mathrm{km}$ between two cities can be treated as a point mass if we only care about its change in position during the trip.}{This is a valid use of the point-particle approximation because the size of the train is negligible compared with the travel distance. In kinematics, when $L_{\text{object}}\ll L_{\text{motion}}$, we can model the object by its center of mass and use a single position vector $\vec r$ to describe the motion.}
\Wrong{We should never model motion in a vacuum, even when we intentionally want to neglect air resistance.}{This confuses a physical model with literal reality; assuming a vacuum is a standard idealization used precisely to remove drag from the analysis.}
\Wrong{None of the above statements is appropriate in physics.}{This is incorrect because option B is a standard and appropriate approximation used in introductory mechanics.}
}{B}

\QBlock{Q2}{Which of the following unit equations is correct?}{%
\Wrong{$1\,\mathrm{mg}=10^{-3}\,\mathrm{kg}$}{This is off by a factor of $1000$; a milligram is much smaller than a gram, so converting directly to kilograms requires $10^{-6}$, not $10^{-3}$.}
\Correct{$1\,\mathrm{g}=10^{-3}\,\mathrm{kg}$}{The SI prefix relation is $1\,\mathrm{kg}=1000\,\mathrm{g}$, so dividing both sides by $1000$ gives $1\,\mathrm{g}=10^{-3}\,\mathrm{kg}$. This option has both the correct numerical factor and consistent mass units.}
\Wrong{$1000\,\mathrm{g}=10^{-6}\,\mathrm{kg}$}{This reverses the metric conversion badly; $1000\,\mathrm{g}$ is actually $1\,\mathrm{kg}$, not a millionth of a kilogram.}
\Wrong{$1\,\mathrm{kg}=100\,\mathrm{g}$}{This uses the wrong power of ten and treats kilo as $10^2$ instead of $10^3$.}
}{B}

\QBlock{Q3}{Which option lists an incorrect SI unit for the quantity named?}{%
\Wrong{Momentum: $\mathrm{kg\cdot m/s}$}{This is a correct SI unit for linear momentum, so selecting it would confuse momentum with some other quantity.}
\Wrong{Pressure: $\mathrm{N/m^2}$}{This is correct because $1\,\mathrm{Pa}=1\,\mathrm{N/m^2}$; marking it wrong would ignore the definition of the pascal.}
\Correct{Power: $\mathrm{kg\cdot m/s^2}$}{Power is the rate of energy transfer, so its SI unit is $\mathrm{J/s}=\mathrm{W}$, not $\mathrm{kg\cdot m/s^2}$. The listed unit is actually a newton, since $1\,\mathrm N=1\,\mathrm{kg\cdot m/s^2}$, so this option mismatches the physical quantity.}
\Wrong{Energy: $\mathrm{N\cdot m}$}{This is dimensionally correct because $1\,\mathrm J=1\,\mathrm{N\cdot m}$, so it is not the incorrect choice.}
}{C}

\QBlock{Q4}{A rescue beacon is located at position $\vec r = -5\hat i + 12\hat j$, where distances are measured in kilometers. Which statement is correct?}{%
\Wrong{The beacon is $7$ km from the origin.}{This comes from subtracting or otherwise combining components incorrectly instead of finding the vector magnitude.}
\Correct{The beacon is $13$ km from the origin.}{The distance from the origin is the magnitude of the position vector, $r=\sqrt{(-5)^2+12^2}$. Evaluating gives $r=\sqrt{25+144}=\sqrt{169}=13\,\mathrm{km}$, so this statement is correct.}
\Wrong{The beacon is $17$ km from the origin.}{This adds the component magnitudes, $5+12$, rather than using the Pythagorean relation for the magnitude of a two-dimensional vector.}
\Wrong{The beacon is $12$ km from the origin.}{This mistakes the $y$-component for the full distance from the origin.}
}{B}

\QBlock{Q5}{If the displacement of an object is given by $\vec s = \left[(0.80\,\mathrm{m/s^3})t^3-(1.60\,\mathrm{m/s^2})t^2\right]\hat i$, which statement is correct?}{%
\Wrong{This motion has constant acceleration.}{Differentiating twice produces an acceleration that depends on $t$, so this choice incorrectly treats a cubic position function as constant-acceleration motion.}
\Correct{The magnitude of the acceleration at $t=0$ is $3.2\,\mathrm{m/s^2}$.}{From $s_x=0.80t^3-1.60t^2$, the velocity is $v_x=2.40t^2-3.20t$ and the acceleration is $a_x=4.80t-3.20$. At $t=0$, $a_x=-3.20\,\mathrm{m/s^2}$, so the magnitude is $3.2\,\mathrm{m/s^2}$.}
\Wrong{The magnitude of the acceleration is always $1.6\,\mathrm{m/s^2}$.}{This uses the coefficient of the $t^2$ term directly and skips the two derivatives needed to obtain acceleration.}
\Wrong{The magnitude of the acceleration at $t=0$ is $0\,\mathrm{m/s^2}$.}{This likely comes from overlooking the constant term that appears after differentiating the displacement function twice.}
}{B}

\QBlock{Q6}{An object takes $4.0$ s to move from position $\vec r_A = 1\hat i - 7\hat j$ to position $\vec r_B = 9\hat i - 1\hat j$ along a straight path, where distances are measured in meters. Which statement is correct?}{%
\Wrong{The moving distance is $6.0\,\mathrm m$.}{This uses only the change in the $y$-coordinate and ignores the simultaneous change in the $x$-coordinate.}
\Wrong{The speed at $t=0$ must be $2.5\,\mathrm{m/s}$.}{This treats average speed, distance divided by time, as if it were guaranteed to equal the initial speed; the initial speed is not determined by the information given.}
\Correct{The moving distance is $10.0\,\mathrm m$.}{The displacement is $\Delta \vec r=(9-1)\hat i+[(-1)-(-7)]\hat j=8\hat i+6\hat j$. Because the path is straight from A to B, the distance traveled is $\lvert\Delta\vec r\rvert=\sqrt{8^2+6^2}=10.0\,\mathrm m$.}
\Wrong{The acceleration at $t=0$ must be $0.63\,\mathrm{m/s^2}$.}{This invents a specific acceleration from distance and time alone, but many different velocity-time histories could connect the same two positions in $4.0$ s.}
}{C}

\QBlock{Q7}{A soccer ball is kicked with initial speed $v_0$ at an angle $\theta$ above the horizontal. Neglecting air resistance, what is the speed of the ball at its maximum height?}{%
\Wrong{$0$}{This confuses zero vertical velocity at the top with zero total velocity; the ball still has horizontal motion.}
\Wrong{$v_0\sin\theta$}{This is the initial vertical component of velocity, not the speed at the top of the trajectory.}
\Wrong{$v_0$}{This ignores that gravity changes the vertical component of the velocity during flight, so the speed is not generally unchanged at the top.}
\Correct{$v_0\cos\theta$}{At maximum height, the vertical component of velocity is zero, $v_y=0$, while the horizontal component remains constant if air resistance is neglected. Therefore the speed there equals $\lvert v_x\rvert=v_0\cos\theta$.}
}{D}

\QBlock{Q8}{A drone is moving with initial velocity $\vec v_0 = (15.0\,\mathrm{m/s})\hat j$ when its thrusters turn on, producing an acceleration $\vec a = (-1.5\,\mathrm{m/s^2})\hat i + (2.0\,\mathrm{m/s^2})\hat j$. Which statement is correct?}{%
\Wrong{The $x$-component of the drone's velocity remains zero.}{Because $a_x=-1.5\,\mathrm{m/s^2}\neq 0$, the $x$-component of velocity begins changing immediately from its initial value.}
\Wrong{The $y$-component of the drone's velocity remains constant.}{A nonzero $y$-acceleration means $v_y$ changes with time, so this choice ignores the effect of $a_y=2.0\,\mathrm{m/s^2}$.}
\Correct{Both the direction of motion and the speed change.}{The velocity changes according to $\vec v(t)=\vec v_0+\vec a t$. Since $\vec a$ has a component perpendicular to the initial motion, the direction changes, and since it also has a component along the motion, the magnitude of $\vec v$ changes as well.}
\Wrong{Neither the speed nor the direction changes.}{This would require zero acceleration, but the thrusters provide nonzero acceleration in both coordinate directions.}
}{C}

\QBlock{Q9}{If the magnitudes of both charges are doubled and their separation is cut in half, what happens to the magnitude of the electrostatic force between them?}{%
\Wrong{It stays the same.}{This ignores that both the numerator and denominator of Coulomb's law change, and in this case they do not cancel.}
\Wrong{It decreases by a factor of $4$.}{This treats the increased charge as if it weakened the force or mishandles the inverse-square dependence on distance.}
\Wrong{It increases by a factor of $8$.}{This captures only part of the scaling and misses that halving $r$ changes $r^2$ by a factor of $1/4$.}
\Correct{It increases by a factor of $16$.}{Coulomb's law gives $F=k\dfrac{|q_1q_2|}{r^2}$. Doubling both charges multiplies the numerator by $4$, and halving the separation makes the denominator $\left(\tfrac{1}{2}r\right)^2=\tfrac{1}{4}r^2$, which multiplies the force by another factor of $4$, so the net change is $4\times 4=16$.}
}{D}

\QBlock{Q10}{An object can escape a planet of mass $M$ and radius $R$ if its speed is at least $v_e=\sqrt{2GM/R}$. Based on this expression, which statement is correct?}{%
\Wrong{A $500\,\mathrm{kg}$ probe needs a higher escape speed than a $5\,\mathrm{kg}$ probe launched from the same planet.}{This confuses the mass of the escaping object with the planet's mass; the formula contains $M$ and $R$, not the probe mass.}
\Correct{For a fixed planet mass, increasing the planet's radius lowers the escape speed.}{From $v_e=\sqrt{2GM/R}$, we have $v_e\propto 1/\sqrt{R}$ when $M$ is constant. So making the planet larger in radius decreases the required escape speed because the same mass is spread farther from the launch point.}
\Wrong{Decreasing the planet's mass raises the escape speed.}{This reverses the dependence on $M$; since $v_e\propto \sqrt{M}$, a smaller planet mass leads to a smaller escape speed.}
\Wrong{A planet with greater mass and greater radius must always have a lower escape speed.}{This claims a definite outcome even though increasing $M$ and increasing $R$ have competing effects, so the result depends on their relative changes.}
}{B}

\QBlock{Q11}{A constant force $\vec F=(6.0\,\mathrm N)\hat i+(-8.0\,\mathrm N)\hat j$ acts on a particle as it moves along the $y$-axis from $y=-1.5\,\mathrm m$ to $y=2.0\,\mathrm m$. What is the work done by the force?}{%
\Correct{$-28.0\,\mathrm J$}{Work is the dot product $W=\vec F\cdot \Delta \vec r$. Here the displacement is purely vertical, $\Delta \vec r=(0)\hat i+(3.5\,\mathrm m)\hat j$, so $W=F_y\Delta y=(-8.0)(3.5)=-28.0\,\mathrm J$.}
\Wrong{$21.0\,\mathrm J$}{This uses the $x$-component of the force even though there is no $x$-displacement, so that component does no work.}
\Wrong{$35.0\,\mathrm J$}{This is close to multiplying the force magnitude by the distance, which is incorrect because work depends on the component parallel to the displacement.}
\Wrong{$6.0\,\mathrm J$}{This likely takes only the $x$-component or mishandles the displacement interval and does not apply the dot-product definition of work.}
}{A}

\QBlock{Q12}{A $0.80\,\mathrm{kg}$ rubber ball is moving at $12\,\mathrm{m/s}$ toward a wall. It rebounds straight back at $9.0\,\mathrm{m/s}$, and the contact lasts $0.050\,\mathrm s$. What is the magnitude of the average force on the ball during the collision?}{%
\Wrong{$48\,\mathrm N$}{This uses an incomplete change in velocity and ignores that the rebound reverses the sign of the final velocity.}
\Correct{$336\,\mathrm N$}{Take motion toward the wall as positive, so $v_i=+12\,\mathrm{m/s}$ and $v_f=-9.0\,\mathrm{m/s}$. Then $F_{\text{avg}}=m\Delta v/\Delta t=0.80\,(-21)/0.050=-336\,\mathrm N$, so the magnitude of the average force is $336\,\mathrm N$.}
\Wrong{$168\,\mathrm N$}{This commonly comes from using only one of the speeds or forgetting that reversing direction makes the velocity change larger than $12-9$.}
\Wrong{$16.8\,\mathrm N$}{This underestimates the impulse dramatically and likely results from dividing by time incorrectly.}
}{B}

\QBlock{Q13}{A $62\,\mathrm{kg}$ skateboarder is initially at rest. She catches a $2.5\,\mathrm{kg}$ medicine ball thrown horizontally at $8.0\,\mathrm{m/s}$. What is her speed immediately after catching the ball?}{%
\Wrong{$0.12\,\mathrm{m/s}$}{This is too small because it underestimates the incoming momentum or overestimates the combined mass effect.}
\Correct{$0.31\,\mathrm{m/s}$}{Because the ball is caught, this is a perfectly inelastic collision and linear momentum is conserved: $m_b v_b=(M+m_b)v_f$. Substituting $m_b=2.5\,\mathrm{kg}$, $v_b=8.0\,\mathrm{m/s}$, and $M=62\,\mathrm{kg}$ gives $v_f=\dfrac{2.5\times 8.0}{64.5}\approx 0.31\,\mathrm{m/s}$.}
\Wrong{$8.0\,\mathrm{m/s}$}{This assumes the skateboarder moves off with the ball's original speed, ignoring conservation of momentum and the much larger total mass after the catch.}
\Wrong{$3.2\,\mathrm{m/s}$}{This is far too large for a light projectile being caught by a much more massive person and reflects a momentum-calculation error.}
}{B}

\end{document}
